\section{Conclusiones}

\begin{enumerate}
  \item Se aplicaron los procedimientos PC-012, PC-014 y PC-029; se obtuvieron los
        errores de indicación y las incertidumbres de los tres instrumentos.
  \item El pie de rey cumple el EMPE de $\pm 80\ \mu\text{m}$, con un punto
        (150 mm) en la zona de incertidumbre. El comparador no cumple: su error
        máximo de $20\ \mu\text{m}$ iguala el EMPE de $\pm 20\ \mu\text{m}$ y
        $U(E) = 7.11\ \mu\text{m}$ lo sitúa fuera de la zona de conformidad.
  \item La incertidumbre del pie de rey es $51.68\ \mu\text{m}$ y está dominada por
        el contacto y la resolución; los términos térmicos son despreciables.
  \item La incertidumbre del comparador es $7.11\ \mu\text{m}$ y está dominada por
        la repetibilidad; su dependencia con la longitud es débil.
  \item La incertidumbre del ultrasonido varía de $126.1$ a $130.2\ \mu\text{m}$;
        el $93.3$\,\% proviene de la deriva del patrón, por lo que no se declara
        conformidad al no existir un EMPE para este equipo.
  \item Se describieron los cuatro métodos de END, sus fundamentos, equipos y
        limitaciones, y se identificó su aplicación en la inspección de componentes.
\end{enumerate}
